
\documentclass{article}
\usepackage{colortbl}
\usepackage{makecell}
\usepackage{multirow}
\usepackage{supertabular}

\begin{document}

\newcounter{utterance}

\centering \large Interaction Transcript for game `codenames', experiment `opponent\_difficult', episode 7 with reflexion{-}llama{-}3.1{-}8b{-}instruct.
\vspace{24pt}

{ \footnotesize  \setcounter{utterance}{1}
\setlength{\tabcolsep}{0pt}
\begin{supertabular}{c@{$\;$}|p{.15\linewidth}@{}p{.15\linewidth}p{.15\linewidth}p{.15\linewidth}p{.15\linewidth}p{.15\linewidth}}
    \# & $\;$A & \multicolumn{4}{c}{Game Master} & $\;\:$B\\
    \hline

    \theutterance \stepcounter{utterance}  
    & & & \multicolumn{2}{p{0.3\linewidth}}{
        \cellcolor[rgb]{0.9,0.9,0.9}{
            \makecell[{{p{\linewidth}}}]{
                \texttt{\tiny{[GM$|$GM]}}
                {'hidden': {'team': ['code'], 'innocent': ['millionaire', 'pole', 'chair', 'wake', 'satellite', 'cricket'], 'opponent': [], 'assassin': ['pupil']}, 'revealed': {'team': {'team': [], 'innocent': [], 'opponent': [], 'assassin': []}, 'opponent': {'team': [], 'innocent': [], 'opponent': [], 'assassin': []}}}
            }
        }
    }
    & & \\ \\

    \theutterance \stepcounter{utterance}  
    & & \multicolumn{4}{p{0.6\linewidth}}{
        \cellcolor[rgb]{0.9,0.9,0.9}{
            \makecell[{{p{\linewidth}}}]{
                \texttt{\tiny{[P1$\langle$GM]}}
                \texttt{You have attempted this same task before and did not succeed. These are the plans you wrote after those attempts.} \\
\texttt{Trial \#1: New plan: } \\
\\ 
\texttt{In my previous attempts, I focused on providing a single{-}word clue that was semantically related to multiple target words. However, I failed to consider the potential for my teammate to guess an opponent or distractor word, which would end their turn and give the opposing team a chance to make a turn. } \\
\\ 
\texttt{Next time, I will make sure to provide a clue that is more specific and targeted towards a smaller set of words, while still being semantically related to them. I will also pay closer attention to the remaining words on the lists and try to avoid providing a clue that could lead to an opponent or distractor word being guessed. For example, I could provide a clue that is related to a specific category or theme that is more likely to be associated with my team's words, rather than a general concept that could be related to multiple words.} \\
\\ 
\texttt{Trial \#2: New plan: } \\
\\ 
\texttt{In my previous attempts, I focused on providing a single{-}word clue that was semantically related to multiple target words. However, I failed to consider the potential for my teammate to guess an opponent or distractor word, which would end their turn and give the opposing team a chance to make a turn. } \\
\\ 
\texttt{Next time, I will make sure to provide a clue that is more specific and targeted towards a smaller set of words, while still being semantically related to them. I will also pay closer attention to the remaining words on the lists and try to avoid providing a clue that could lead to an opponent or distractor word being guessed. For example, I could provide a clue that is related to a specific category or theme that is more likely to be associated with my team's words, rather than a general concept that could be related to multiple words.} \\
\\ 
\texttt{To achieve this, I will start by analyzing the words on my team's list and identifying common themes or categories that they belong to. I will then try to come up with a clue that is related to this category, but not so general that it could be applied to the opponent's words. For example, if my team's words are all related to cryptography, I could provide a clue that is related to encryption or decryption, but not so broad that it could be applied to words like "wake" or "satellite".} \\
\\ 
\\ 
\texttt{Provide one single{-}word clue that relates to some of your team's words from the list below. You can choose to how many and to which words your clue relates to.} \\
\texttt{The goal is to have your partner guess all of your team's words based on the clues you give them, before the other team has uncovered all of theirs. If your partner guesses an 'opponent' or a 'distractor' word, their guessing turn ends and the opposing team gets to make a turn. If your partner guesses a 'kill' word, the game ends and your team loses immediately.} \\
\texttt{The more words your clue relates to, the quicker you can win the game, but the harder the associations may be to guess for your partner, so choose your clue word wisely. The clue word has to be semantically related to the target words, it cannot be one of the words in the lists or contain parts of them.} \\
\\ 
\texttt{Always give your single{-}word clue and your comma{-}separated list of related target words in the following format and make your answers as short as possible, never include any other text than is required in this form:} \\
\\ 
\texttt{CLUE: <WORD>} \\
\texttt{TARGETS: <TARGETS>} \\
\\ 
\texttt{Your team words are: code.} \\
\texttt{Your opponent's team words are: .} \\
\texttt{Distractor words are: millionaire, pole, chair, wake, satellite, cricket.} \\
\texttt{The kill words are: pupil.} \\
            }
        }
    }
    & \\ \\

    \theutterance \stepcounter{utterance}  
    & \multicolumn{4}{p{0.6\linewidth}}{
        \cellcolor[rgb]{0.9,0.9,0.9}{
            \makecell[{{p{\linewidth}}}]{
                \texttt{\tiny{[P1$\rangle$GM]}}
                \texttt{CLUE: cipher} \\
\texttt{TARGETS: code} \\
            }
        }
    }
    & & \\ \\

    \theutterance \stepcounter{utterance}  
    & & & \multicolumn{2}{p{0.3\linewidth}}{
        \cellcolor[rgb]{0.9,0.9,0.9}{
            \makecell[{{p{\linewidth}}}]{
                \texttt{\tiny{[GM$|$GM]}}
                \texttt{cipher} \\
            }
        }
    }
    & & \\ \\

    \theutterance \stepcounter{utterance}  
    & & & \multicolumn{2}{p{0.3\linewidth}}{
        \cellcolor[rgb]{0.9,0.9,0.9}{
            \makecell[{{p{\linewidth}}}]{
                \texttt{\tiny{[GM$|$GM]}}
                ['code']
            }
        }
    }
    & & \\ \\

    \theutterance \stepcounter{utterance}  
    & & & \multicolumn{2}{p{0.3\linewidth}}{
        \cellcolor[rgb]{0.9,0.9,0.9}{
            \makecell[{{p{\linewidth}}}]{
                \texttt{\tiny{[GM$|$GM]}}
                {'word': 'code', 'assignment': 'team'}
            }
        }
    }
    & & \\ \\

    \theutterance \stepcounter{utterance}  
    & & & \multicolumn{2}{p{0.3\linewidth}}{
        \cellcolor[rgb]{0.9,0.9,0.9}{
            \makecell[{{p{\linewidth}}}]{
                \texttt{\tiny{[GM$|$GM]}}
                \texttt{opponent has won} \\
            }
        }
    }
    & & \\ \\

    \theutterance \stepcounter{utterance}  
    & \multicolumn{4}{p{0.6\linewidth}}{
        \cellcolor[rgb]{0.95,0.95,0.95}{
            \makecell[{{p{\linewidth}}}]{
                \texttt{\tiny{[P2$\rightarrow$P1]}}
                \texttt{Post{-}episode Reflexion (not sent to GM):} \\
\texttt{New plan: } \\
\\ 
\texttt{In my previous attempts, I focused on providing a single{-}word clue that was semantically related to multiple target words. However, I failed to consider the potential for my teammate to guess an opponent or distractor word, which would end their turn and give the opposing team a chance to make a turn. } \\
\\ 
\texttt{Next time, I will make sure to provide a clue that is more specific and targeted towards a smaller set of words, while still being semantically related to them. I will also pay closer attention to the remaining words on the lists and try to avoid providing a clue that could lead to an opponent or distractor word being guessed. } \\
\\ 
\texttt{To achieve this, I will start by analyzing the words on my team's list and identifying common themes or categories that they belong to. I will then try to come up with a clue that is related to this category, but not so general that it could be applied to the opponent's words. } \\
\\ 
\texttt{In this case, my team's word is "code", which is related to cryptography. I will provide a clue that is specific to this category, such as "algorithm" or "encryption". This will increase the chances of my teammate guessing the correct word, while minimizing the risk of them guessing an opponent or distractor word.} \\
            }
        }
    }
    & & \\ \\

\end{supertabular}
}

\end{document}
